\NormalFont\ShortTitle{Gálatas}
{\MT Carta de Paulo aos Gálatas

\par }\ChapOne{1}{\PP \VerseOne{1}Paulo, apóstolo (não da parte dos seres humanos, nem por meio de homem algum, mas sim por Jesus Cristo, e por Deus Pai, que o ressuscitou dos mortos),
\VS{2}e todos os irmãos que estão comigo, para as igrejas da Galácia.
\VS{3}Graça seja convosco, e também a paz de Deus Pai, e do nosso Senhor Jesus Cristo,
\VS{4}que deu a si mesmo pelos nossos pecados, para nos tirar da presente era maligna, conforme a vontade do nosso Deus e Pai,
\VS{5}ao qual seja a glória para todo o sempre, Amém!
\VS{6}Admiro-me de que tão rapidamente desviastes daquele que vos chamou na graça de Cristo, em troca de outro evangelho.
\VS{7}Não que haja outro
{\ADD{evangelho}} ,porém há alguns que vos inquietam, e querem perverter o Evangelho de Cristo.
\VS{8}Porém, ainda que nós mesmos, ou um anjo do céu vos anuncie outro evangelho além do que já vos anunciamos, maldito seja.
\VS{9}Como já havíamos dito, volto também agora a dizer: se alguém vos anunciar outro evangelho além do que já recebestes, maldito seja.
\VS{10}Pois agora estou buscando a aprovação das pessoas, ou a de Deus? Ou procuro agradar a pessoas? Pois, se ainda tentasse agradar a pessoas, eu não seria servo de Cristo.
\VS{11}Mas faço-vos saber, irmãos, que o Evangelho anunciado por mim não se baseia em pessoas;
\FTNT{*}{{\FR 1:11 }
Lit. não é conforme pessoas
}
\VS{12}pois não o recebi, nem aprendi de ser humano algum, mas sim pela revelação de Jesus Cristo.
\VS{13}Pois já ouvistes da minha conduta no judaísmo, como em excesso eu perseguia e tentava destruir a Igreja de Deus;
\VS{14}e
{\ADD{como}} no judaísmo eu era mais avançado que muitos de minha idade em minha nação, e era extremamente zeloso das tradições de meus pais.
\VS{15}Mas quando Deus ( que me separou desde o ventre da minha mãe, e por sua graça me chamou ) se agradou
\VS{16}de revelar o seu Filho em mim, para eu evangelizar os gentios, de imediato, não fui pedir conselho com pessoa alguma;
\FTNT{†}{{\FR 1:16 }
pessoa alguma
Lit. “carne e sangue”
}
\VS{17}nem subi a Jerusalém para os que já eram apóstolos antes de mim; em vez disso, parti para a Arábia, e voltei outra vez a Damasco.
\VS{18}Então, depois de três anos, subi a Jerusalém para visitar Pedro, e estive com ele por quinze dias.
\VS{19}E vi nenhum outro dos apóstolos, a não ser Tiago, o irmão do Senhor.
\VS{20}Ora, das coisas que vos escrevo, eis que, diante de Deus, não estou mentindo.
\VS{21}Depois fui para as regiões da Síria e da Cilícia.
\VS{22}Eu, porém, não era conhecido de rosto pelas igrejas da Judeia que estão em Cristo;
\VS{23}mas somente haviam ouvido que “Aquele que antes nos perseguia agora anuncia a fé que antes tentava destruir”.
\VS{24}E glorificavam a Deus por causa de mim.

\par }\Chap{2}{\PP \VerseOne{1}Depois de passados catorze anos, subi outra vez a Jerusalém com Barnabé, levando comigo também Tito.
\VS{2}E, por causa de uma revelação, subi, e lhes declarei o Evangelho que prego entre os gentios.
{\ADD{Isso}} , porém, foi em particular, com os mais influentes,
\FTNT{*}{{\FR 2:2 }
Isto é, os líderes da ireja
} para que eu eu não corresse ou tivesse corrido em vão.
\VS{3}Porém nem mesmo Tito, que estava comigo, sendo grego, foi obrigado a se circuncidar.
\VS{4}E
{\ADD{isso}} por causa de falsos irmãos, que haviam se infiltrado, e entraram secretamente para espionar a nossa liberdade que temos em Cristo Jesus, a fim de nos tornar escravos.
\VS{5}A eles nem sequer por um momento
\FTNT{†}{{\FR 2:5 }
Lit. hora
} cedemos a nos sujeitarmos, para que a verdade do Evangelho permanecesse em vós.
\VS{6}E quanto aos de maior influência (o que haviam sido antes não me importa; Deus não se interessa na aparência humana); esses, pois, que eram influentes, nada me acrescentaram.
\VS{7}Pelo contrário, quando viram que o Evangelho aos incircuncisos
\FTNT{‡}{{\FR 2:7 }
Lit. “Evangelho da incircuncisão”
} havia sido confiado a mim, assim como a Pedro, aos circuncisos,
\FTNT{§}{{\FR 2:7 }
Lit. “o [Evangelho] da circuncisão”
}
\VS{8}(pois aquele que operou em Pedro para o apostolado da circuncisão, esse operou também em mim para os gentios)
\VS{9}e quando Tiago, Cefas, e João, que eram considerados como colunas, reconheceram a graça que me foi dada, eles estenderam as mãos direitas de comunhão a mim e a Barnabé, para que nós
{\ADD{fôssemos}} aos gentios, e eles aos circuncisos;
\VS{10}sob a condição de que nos lembrássemos dos pobres; isso também procurei fazer com empenho.
\VS{11}E quando Pedro veio a Antioquia, estive contra ele face a face, pois ele devia ser repreendido;
\VS{12}porque, antes que alguns que Tiago enviou chegassem, ele comia com os gentios; mas depois que chegaram, ele se retirou e se separou, temendo os que eram da circuncisão.
\VS{13}E também com ele os outros judeus fingiram, de maneira que até Barnabé se deixou levar pela hipocrisia deles.
\VS{14}Mas quando vi que não estavam agindo corretamente conforme a verdade do Evangelho, disse na presença de todos a Pedro: “Se tu, que és judeu, vives como gentio, e não como judeu, por que obrigas os gentios a viverem como judeus?”
\VS{15}Nós, que somos judeus por natureza, e não pecadores dentre os gentios,
\VS{16}sabemos que o ser humano não é justificado pelas obras da Lei, mas sim pela fé em Jesus Cristo.
\FTNT{*}{{\FR 2:16 }
Trad. alt. fidelidade de Jesus Cristo
} Também nós temos crido em Cristo Jesus, para que fôssemos justificados pela fé em Cristo,
\FTNT{†}{{\FR 2:16 }
Trad. alt. fidelidade de Cristo
} e não pelas obras da Lei; pois ninguém
\FTNT{‡}{{\FR 2:16 }
Lit. nenhuma carne
} se justificará pelas obras da Lei.
\VS{17}Mas se nós, que buscamos ser justificados em Cristo, também nós mesmos somos achados pecadores, por isso Cristo contribui com o pecado?
\FTNT{§}{{\FR 2:17 }
Lit. Cristo é servidor do pecado?
} De maneira nenhuma!
\VS{18}Pois, se volto a construir as coisas que já destruí, provo que eu mesmo cometi transgressão.
\VS{19}Pois pela Lei estou morto para a Lei, a fim de que eu viva para Deus.
\VS{20}Já estou crucificado com Cristo. Estou vivendo não mais eu, mas Cristo vive em mim; e vivo a minha vida na carne por meio da fé no Filho de Deus,
\FTNT{*}{{\FR 2:20 }
Trad. alt. fidelidade do Filho de Deus
} que me amou, e entregou a si mesmo por mim.
\VS{21}Não anulo a graça
\FTNT{†}{{\FR 2:21 }
Isto é, o favor não merecido
} de Deus; pois, se a justiça
\FTNT{‡}{{\FR 2:21 }
Isto é, ser justo, que aqui significa ser aceito (absolvido) por Deus
} é por meio da Lei, logo Cristo morreu por nada.

\par }\Chap{3}{\PP \VerseOne{1}Ó insensatos gálatas, quem vos iludiu
\FTNT{*}{{\FR 3:1 }
iludiu
lit. enfeitiçou
} para não obedecerdes à verdade? Diante dos vossos olhos Jesus Cristo foi retratado entre vós como crucificado!
\VS{2}Só isto eu queria saber de vós: recebestes o Espírito pelas obras da Lei ou pela pregação da fé?
\VS{3}Sois vós tão insensatos que, depois de começardes no Espírito, agora terminareis na carne?
\VS{4}Experimentastes
\FTNT{†}{{\FR 3:4 }
Trad. alt. Sofrestes
} tantas coisas em vão? Se é que foi mesmo em vão!
\VS{5}Ora, aquele que vos dá o Espírito e que opera maravilhas entre vós
{\ADD{faz isso}} por causa das obras da Lei ou da pregação da fé?
\VS{6}Assim como “Abraão creu em Deus, e foi lhe reputado como justiça”,
{\RQ{Gênesis 15:6
}}
\VS{7}entendei, pois, que os que são da fé são filhos de Abraão.
\VS{8}E a Escritura, prevendo que Deus justificaria os gentios pela fé, prenunciou o Evangelho a Abraão,
{\ADD{dizendo}} : Todas nas nações serão abençoadas em ti.
{\RQ{Gênesis 12:3; 18:18; 22:18
}}
\VS{9}Portanto, os que são da fé são abençoados com o crente Abraão.
\VS{10}Pois todos os que são das obras da Lei estão sob maldição, porque está escrito: Maldito todo aquele que não permanecer em todas as coisas que estão escritas no livro da Lei, para fazê-las.
{\RQ{Deuteronômio 27:26
}}
\VS{11}E é evidente que pela Lei ninguém será justificado, porque: O justo viverá pela fé.
{\RQ{Habacuque 2:4
}}
\VS{12}A Lei não provém
\FTNT{‡}{{\FR 3:12 }
Lit. é
} da fé; porém: Quem fizer estas coisas por elas viverá.
\FTNT{§}{{\FR 3:12 }
Levítico 18:5
}
\VS{13}Cristo nos resgatou da maldição da Lei ao se fazer maldição para o nosso benefício (pois está escrito: Maldito todo aquele que for pendurado em um madeiro.)
{\RQ{Deuteronômio 21:23
}}
\VS{14}com a finalidade de que a bênção de Abraão chegasse aos gentios em Cristo Jesus, para que recebêssemos a promessa do Espírito
\FTNT{*}{{\FR 3:14 }
Ou: o Espírito prometido
} por meio da fé.
\VS{15}Irmãos, estou falando em termos humanos: ainda que seja o pacto de uma pessoa, depois de confirmado, ninguém o anula, nem lhe acrescenta.
\VS{16}Ora, as promessas foram ditas a Abraão e à sua descendência.
\FTNT{†}{{\FR 3:16 }
Lit. semente
} E não diz: “E às descendências”
\FTNT{‡}{{\FR 3:16 }
Lit. sementes
} como que
{\ADD{falando}} de muitos, mas sim como a um: “E à tua descendência”,
{\RQ{Gênesis 12:7; 13:15; 24:7
}} que é Cristo.
\VS{17}Mas digo isto: o pacto foi confirmado anteriormente por Deus em Cristo; e a Lei que veio quatrocentos e trinta anos depois não o invalida de maneira que anule a promessa.
\VS{18}Pois, se a herança é pela Lei, já não é pela promessa; mas foi por meio da promessa que Deus a concedeu gratuitamente a Abraão.
\VS{19}Para que, pois, é a Lei? Ela foi ordenada por causa das transgressões, até que chegasse o descendente
\FTNT{§}{{\FR 3:19 }
Isto é, a descendência do vers. 16, que é Cristo
} a quem a promessa havia sido feita; e foi posta pelos anjos na mão de um mediador.
\VS{20}Ora, o mediador não é de um só, mas Deus é um.
\VS{21}Acaso, pois, a Lei é contra as promessas de Deus? De maneira nenhuma! Pois, se a Lei houvesse sido entregue para que pudesse dar vida, na verdade a justiça teria sido pela Lei.
\VS{22}Mas a Escritura prendeu tudo debaixo do pecado a fim de que a promessa fosse dada aos crentes por meio da fé em Jesus Cristo.
\FTNT{*}{{\FR 3:22 }
fé em Jesus Cristo
trad. alt. fidelidade de Jesus Cristo
}
\VS{23}Porém, antes que a fé viesse, estávamos vigiados sob a Lei, e presos, até que a fé fosse revelada.
\VS{24}Dessa maneira, a Lei foi nosso tutor em condução a Cristo, para que pela fé fossemos justificados.
\VS{25}Mas depois que a fé chegou, já não estamos mais sob um tutor,
\VS{26}pois todos vós sois filhos de Deus por meio da fé em Cristo Jesus;
\VS{27}pois todos vós que fostes batizados em Cristo já vos revestistes de Cristo.
\VS{28}Assim, não há judeu nem grego; não há servo nem livre; não há macho nem fêmea. Pois todos vós sois um em Cristo Jesus.
\VS{29}E, se vós sois de Cristo, então sois descendência
\FTNT{†}{{\FR 3:29 }
Lit. semente
} de Abraão, e herdeiros conforme a promessa.

\par }\Chap{4}{\PP \VerseOne{1}Digo, porém, que, enquanto o herdeiro é criança, em nada é diferente do servo, ainda que seja senhor de tudo.
\VS{2}Mas está sob tutores e administradores, até o tempo determinado pelo pai.
\VS{3}Assim também nós, quando éramos crianças, estávamos escravizados debaixo dos princípios elementares do mundo.
\VS{4}Mas quando o tempo se completou, Deus enviou o seu Filho, nascido de mulher, nascido sob a Lei,
\VS{5}para redimir os que estavam sob a Lei, a fim de que nós recebêssemos a adoção como filhos.
\VS{6}E, como sois filhos, Deus enviou aos vossos corações o Espírito do seu Filho, que clama: “Aba, Pai!”
\VS{7}Portanto, tu já não és mais servo, mas sim filho; e se és filho, também és herdeiro de Deus por meio de Cristo.
\VS{8}Quando, porém, não conhecíeis a Deus, vós servíeis aos que por natureza não são deuses.
\VS{9}Mas agora, que conheceis Deus, mais que isso, por Deus sois conhecidos, por que voltais outra vez aos fracos e miseráveis princípios elementares, aos quais quereis servir de novo?
\VS{10}Vós guardais dias, meses, tempos, e anos.
\VS{11}Temo a vosso respeito de que talvez eu tenha trabalhado para convosco em vão.
\VS{12}Rogo-vos, irmãos, que sejais como eu, pois também eu sou como vós. Nenhum mal me fizestes.
\VS{13}E vós sabeis que foi com uma enfermidade da carne que vos anunciei o Evangelho pela primeira vez;
\VS{14}E não cedestes à tentação de ter desprezo ou repugnância por meu aspecto físico;
\FTNT{*}{{\FR 4:14 }
meu aspecto físico
Lit. minha carne
} ao contrário, vós me recebestes como a um mensageiro de Deus, como ao
{\ADD{próprio}} Cristo Jesus.
\VS{15}O que houve, pois, com a vossa alegria? Pois eu vos dou testemunho de que, se fosse possível, teríeis arrancado os vossos próprios olhos para dá-los a mim.
\VS{16}Tornei-me, por acaso, vosso inimigo por ter vos dito a verdade?
\VS{17}Eles têm interesse por vós, mas não com boa intenção;
\FTNT{†}{{\FR 4:17 }
Lit. não da boa (ou correta) maneira
} mas querem vos isolar
{\ADD{de mim}} ,para que vós mostreis interesse por eles.
\VS{18}É bom, porém, mostrar interesse
\FTNT{‡}{{\FR 4:18 }
Ou: ser objeto de interesse
} para o bem sempre, e não somente quando estou presente convosco.
\VS{19}Meus filhinhos, por quem volto a sofrer dores de parto, até que Cristo seja formado em vós:
\VS{20}queria eu estar agora presente convosco, e mudar meu tom de voz, pois estou perplexo a vosso respeito.
\VS{21}Respondei-me vós, que quereis estar sob a Lei: não ouvis a Lei?
\VS{22}Pois está escrito que Abraão tinha dois filhos: um da escrava, e outro da livre.
\VS{23}O da escrava nasceu segundo a carne, mas o da livre, pela promessa.
\VS{24}Isto serve de ilustração, pois estes são dois pactos: um é o do monte Sinai, que gera filhos para a escravidão, que é Agar.
\VS{25}Pois esta Agar é o monte Sinai, na Arábia, e corresponde à Jerusalém de agora, e é escravizada com os filhos dela.
\VS{26}Mas a Jerusalém de cima é livre; esta é a mãe de todos nós.
\VS{27}Pois está escrito: Alegra-te, ó estéril, que não geras filhos; extravasa gritando de júbilo, tu, que não estás de parto; porque muitos mais são os filhos da solitária, que os da que tem marido.
{\RQ{Isaías 54:1
}}
\VS{28}Nós, porém, irmãos, como Isaque, somos filhos da promessa.
\VS{29}Porém, assim como naquele tempo o que nasceu segundo a carne perseguia o que
{\ADD{nasceu}} segundo o Espírito, assim também é agora.
\VS{30}Mas o que diz a Escritura? Expulsa a escrava e o filho dela, porque de maneira nenhuma o filho da escrava herdará com o filho da livre.
{\RQ{Gênesis 21:10
}}
\VS{31}Portanto, irmãos, não somos filhos da escrava, mas sim da livre.

\par }\Chap{5}{\PP \VerseOne{1}Estai, pois, firmes na liberdade com que Cristo nos libertou, e não volteis a vos prender com o jugo da escravidão.
\VS{2}Eis que eu, Paulo, vos digo, que se vos deixardes circuncidar, Cristo vos será útil em nada.
\VS{3}E de novo atesto que todo homem que se deixar circuncidar está obrigado a obedecer a toda a Lei.
\VS{4}Desligados estais de Cristo, vós que
{\ADD{quereis}} ser justos pela Lei; da graça caístes.
\VS{5}Pois, por meio do Espírito, pela fé, aguardamos a esperança da justiça;
\VS{6}porque, em Cristo Jesus, nem a circuncisão nem a incircuncisão tem valor algum; mas sim a fé, que opera por meio do amor.
\VS{7}Estáveis correndo bem; quem vos impediu de obedecerdes à verdade?
\VS{8}Esta persuasão não parte daquele que vos chama.
\VS{9}Um pouco de fermento leveda toda a massa.
\VS{10}Acerca de vós, confio no Senhor de que não mudareis a vossa mentalidade; mas aquele que vos perturba, seja quem for, sofrerá o julgamento.
\VS{11}Mas se eu, irmãos, ainda prego a circuncisão, por que, então, sou perseguido? Então a ofensa da cruz está anulada!
\VS{12}Gostaria que aqueles que estão vos perturbando castrassem a si mesmos.
\VS{13}Pois vós, irmãos, fostes chamados para a liberdade. Somente não
{\ADD{useis}} a liberdade como oportunidade para a carne; em vez disso, servi-vos uns aos outros pelo amor.
\VS{14}Pois toda a Lei se cumpre em uma só regra, que é: Amarás ao teu próximo como a ti mesmo.
{\RQ{Levítico 19:18
}}
\VS{15}Se, porém, mordeis e devorais uns aos outros, cuidado para não vos destruirdes mutuamente.
\VS{16}Mas eu digo: andai no Espírito, e não executeis o mau desejo da carne.
\VS{17}Pois a carne deseja contra o Espírito, e o Espírito contra a carne; e estes se opõem mutuamente, para que não façais o que quereis.
\VS{18}Mas, se sois guiados pelo Espírito, não estais debaixo da Lei.
\VS{19}As obras da carne são evidentes. São elas: adultério, pecado sexual, impureza, devassidão,
\VS{20}idolatria, feitiçaria, inimizades, brigas, ciúmes, iras, rivalidades egoístas, desavenças, facções,
\VS{21}invejas, homicídios, bebedices, orgias, e coisas semelhantes a essas, das quais eu havia vos dito anteriormente; assim como eu também haviavos dito antes que os que praticam tais coisas não herdarão o Reino de Deus.
\VS{22}Mas o fruto do Espírito é: amor, alegria, paz, paciência, benignidade, bondade, fidelidade,
\FTNT{*}{{\FR 5:22 }
Ou: fé
}
\VS{23}mansidão, domínio próprio. Contra essas coisas não há lei.
\VS{24}Os que são de Cristo crucificaram a carne com as paixões e os maus desejos.
\VS{25}Se vivemos no Espírito, também no Espírito andemos.
\VS{26}Não nos tornemos presunçosos, irritando uns aos outros, e invejando uns aos outros.

\par }\Chap{6}{\PP \VerseOne{1}Irmãos, se alguém for surpreendido em algum pecado, vós que sois espirituais, restaurai ao tal com espírito de mansidão, prestando atenção a ti mesmo, para que não sejas tentado também.
\VS{2}Levai as cargas uns dos outros, e assim cumprireis a lei de Cristo.
\VS{3}Pois, se alguém pensa ser alguma coisa, sendo nada, engana a si mesmo.
\VS{4}Cada um, porém, examine a sua própria obra, e então terá orgulho em si mesmo sozinho, e não em outro;
\VS{5}pois cada um levará a sua própria carga.
\VS{6}Mas aquele que é instruído na palavra compartilhe todas as boas coisas com aquele que o instrui.
\VS{7}Não vos enganeis: Deus não se deixa escarnecer; pois, o que o ser humano semear, isso também colherá.
\VS{8}Pois quem semear para a sua carne, da carne colherá degradação; mas quem semear para o Espírito, do Espírito colherá a vida eterna.
\VS{9}Porém não cansemos de fazer o bem, pois no tempo devido colheremos, se não desistirmos.
\FTNT{*}{{\FR 6:9 }
na língua original, desistir por esgotamento
}
\VS{10}Portanto, enquanto temos oportunidade, façamos o bem a todos; mas principalmente aos familiares
\FTNT{†}{{\FR 6:10 }
Ou domésticos, isto é, os da mesma casa
} da fé.
\VS{11}Olhai como são grandes as letras que vos escrevi da minha mão.
\VS{12}Todos os que querem mostrar boa aparência na carne, estes vos obrigam a vos circuncidarem, somente para não serem perseguidos por causa da cruz de Cristo.
\VS{13}Pois nem mesmo os que se circuncidam guardam a Lei; mas querem que vos circuncideis, para se orgulharem de vossa carne.
\VS{14}Mas longe de mim esteja me orgulhar, a não ser na cruz do nosso Senhor Jesus Cristo; por meio dela, o mundo está crucificado para mim, e eu para o mundo.
\VS{15}Pois em Cristo Jesus, nem a circuncisão tem valor algum, nem a incircuncisão; mas sim a nova criatura.
\VS{16}E todos os que andarem conforme esta regra, paz e misericórdia estejam sobre eles, e sobre o Israel de Deus.
\VS{17}De agora em diante, ninguém me perturbe, porque trago no meu corpo as marcas do Senhor Jesus.
\VS{18}A graça do nosso Senhor Jesus Cristo seja com o vosso espírito, irmãos. Amém!{\ADD{Escrita de Roma para os Gálatas}}
\par }